\section[Candidatos fraccionarios no Chua]{Candidatos fraccionarios no Chua: álgebra verificada y alcance dinámico}
\label{part:nonchua}

\section{Resultado de la búsqueda y alcance real}

La revisión cruzada de las fuentes primarias, la extracción de datos publicados
y la matriz de cobertura del repositorio no produce una lista amplia de
candidatos ejecutados. Solamente aparecen dos sistemas fraccionarios no Chua
ligados de forma explícita a un reclamo publicado de atractor oculto:
el Lorenz generalizado y el sistema de Rabinovich--Fabrikant (RF), ambos tomados
de Danca~\cite{Danca2017}.

\begin{table}[ht]
\centering
\small
\caption{Clasificación de la evidencia localizada.}
\begin{tabularx}{\linewidth}{>{\raggedright\arraybackslash}p{0.23\linewidth}
>{\raggedright\arraybackslash}p{0.19\linewidth}X}
\toprule
Caso & Estado local & Evidencia disponible y límite\\
\midrule
Lorenz generalizado, \(q=0.995\) & \statusref &
Ecuaciones, parámetros, equilibrios redondeados, espectros publicados y un
ensayo bibliográfico de vecindad. No existe sistema integrado, semilla, corrida
de continuación ni artefacto de cuenca local.\\
Rabinovich--Fabrikant, \(q=0.998\) & \statusref &
Ecuaciones, parámetros, cinco equilibrios, espectros publicados y radios de
vecindad. No existe reproducción local del atractor ni de su cuenca.\\
Lorenz estándar, RF \(b=0.98\), jerk, financiero y four-wing & diagnóstico &
Son bancos de Lyapunov o comparaciones parciales. No certifican búsqueda de
atractores ocultos y no deben mezclarse con los dos casos de Danca.\\
Genesio--Tesi, jerk hiperbólico y otras familias del Artículo 03 & diseño &
Poseen realizaciones algebraicas propuestas, pero el prerregistro no Chua sigue
sin ejecuciones L0--L6 completas.\\
\bottomrule
\end{tabularx}
\end{table}

Los diseños locales para jerk, Genesio--Tesi, sistemas memristivos, Lorenz, RF
y el sistema financiero, así como los bancos DK2018 de Lyapunov para Lorenz
estándar y RF, no reúnen simultáneamente ecuaciones, parámetros, orden,
condición inicial o semilla, continuación y evidencia de cuenca para un
candidato no Chua. Por ello permanecen fuera de los dos casos principales.

El propio reporte unificado declara que Lorenz es un ejemplo pedagógico y que no
pertenece a su metodología activa. La fuente local decisiva es
\path{version_2/validation/references/published_validation_data_extraction_v1.json};
la clasificación conservadora procede de
\path{version_2/validation/published_reference_coverage.json}. En esta última,
los dos casos están marcados como \texttt{reference\_data\_only}, sin sistema
implementado, pruebas ni artefactos de validación. Por ello, en lo que sigue la
palabra \emph{candidato} significa \emph{caso prioritario para implementar y
buscar}, no \emph{atractor localmente reproducido}.

\begin{hallazgo}{Conclusión de la auditoría de fuentes}
El recuerdo del usuario era parcialmente correcto: había sistemas no Chua
registrados alrededor del reporte y de los manuscritos, pero el reporte
unificado no ejecutó su búsqueda. Los únicos dos con procedencia bibliográfica
directa de atractor oculto son Lorenz generalizado y RF; su cobertura local es
todavía sólo documental.
\end{hallazgo}

\section{Marco común de realización y localización}

\subsection{Dinámica de Caputo y transferencia fraccionaria}

Para \(0<q<1\), la derivada de Caputo con terminal inferior \(t_0\) se escribe
\begin{equation}
 {}^{\mathrm C}D_{t_0}^{q}x(t)
 =\frac{1}{\Gamma(1-q)}
 \int_{t_0}^{t}(t-\tau)^{-q}x'(\tau)\,d\tau .
 \label{nc:eq:caputo}
\end{equation}
La ecuación autónoma
\begin{equation}
 {}^{\mathrm C}D_{t_0}^{q}x(t)=F(x(t)),\qquad x(t_0)=x_0,
 \label{nc:eq:ivp}
\end{equation}
es equivalente a la ecuación integral de Volterra
\begin{equation}
 x(t)=x_0+\frac{1}{\Gamma(q)}
 \int_{t_0}^{t}(t-\tau)^{q-1}F(x(\tau))\,d\tau .
 \label{nc:eq:volterra}
\end{equation}
La historia desde \(t_0\) es parte del problema; no puede reiniciarse durante
una continuación sin declarar que se ha definido un problema diferente.

Para el bloque lineal
\begin{equation}
 {}^{\mathrm C}D_{t_0}^{q}x=Ax+Bu,
 \qquad y=Cx,
 \label{nc:eq:linear-block}
\end{equation}
la transformada de Laplace cumple
\begin{equation}
 \mathcal L\{{}^{\mathrm C}D_{t_0}^{q}x\}(s)
 =s^qX(s)-s^{q-1}x(t_0).
\end{equation}
Con estado inicial nulo, la transferencia entrada--salida es
\begin{equation}
 W_q(s)=C(s^qI-A)^{-1}B.
 \label{nc:eq:transfer-general}
\end{equation}
Se usa la rama principal. En el predictor armónico,
\begin{equation}
 (\ii k\omega)^q=(k\omega)^q
 \exp\!\left(\ii\frac{q\pi}{2}\right),\qquad k>0.
 \label{nc:eq:principal-branch}
\end{equation}
Esta sustitución produce un problema algebraico estacionario; no afirma la
existencia de una órbita periódica exacta del PVI causal de Caputo.

\subsection{Realización de Lur'e multicanal}

Los dos candidatos no son escalares en su no linealidad. Se usa la forma exacta
\begin{equation}
 {}^{\mathrm C}D_{t_0}^{q}x=Ax+B\Phi(Cx),
 \qquad x\in\RR^n,
 \label{nc:eq:mimo-lure}
\end{equation}
con \(\Phi:\RR^p\to\RR^m\) estática. La identidad que debe comprobarse es
\begin{equation}
 Ax+B\Phi(Cx)-F(x)\equiv0.
 \label{nc:eq:lure-residual}
\end{equation}
No basta que las expresiones coincidan en una trayectoria o en un equilibrio.

Sea \(\bar y=C\bar x\) un sesgo y
\begin{equation}
 K=D\Phi(\bar y)\in\RR^{m\times p}.
\end{equation}
La ecuación de cierre linealizada del primer armónico es
\begin{equation}
 \det\!\left[I_p-W_q(\ii\omega)K\right]=0.
 \label{nc:eq:det-closure}
\end{equation}
Para amplitud finita, \(K\) se sustituye por los coeficientes de Fourier
exactos de la no linealidad.

\subsection{Coeficientes exactos y balance multiharmónico}

Se adopta la convención compleja
\begin{equation}
 x(\theta)=\sum_{k=-H}^{H}X_k e^{\ii k\theta},
 \qquad X_{-k}=\overline{X_k},\qquad X_0\in\RR^n,
 \label{nc:eq:fourier-state}
\end{equation}
con \(\theta=\omega t\). Si \(u\) y \(v\) tienen coeficientes \(U_k,V_k\),
la multiplicación no se aproxima: sus coeficientes son la convolución
\begin{equation}
 \widehat{(uv)}_k=\sum_{m\in\mathbb Z}U_mV_{k-m}.
 \label{nc:eq:conv2}
\end{equation}
Para un producto triple,
\begin{equation}
 \widehat{(uvw)}_k
 =\sum_{m\in\mathbb Z}\sum_{n\in\mathbb Z}
 U_mV_nW_{k-m-n}.
 \label{nc:eq:conv3}
\end{equation}
En una truncación se toma \(U_k=0\) para \(|k|>H\). Esto calcula exactamente
los coeficientes retenidos de la señal truncada; la omisión de armónicos altos
se controla aumentando \(H\).

Si \(\widehat\Phi_k\) denota el coeficiente \(k\) de \(\Phi(Cx(\theta))\), el
sistema completo de balance es
\begin{align}
 AX_0+B\widehat\Phi_0&=0,
 \label{nc:eq:hb0}\\
 \left[(\ii k\omega)^qI-A\right]X_k-B\widehat\Phi_k&=0,
 \qquad k=1,\ldots,H.
 \label{nc:eq:hbk}
\end{align}
Las incógnitas son \(X_0,X_1,\ldots,X_H,\omega\), junto con una condición de
fase, por ejemplo
\begin{equation}
 \operatorname{Im}(e_j\trans X_1)=0,\qquad
 \operatorname{Re}(e_j\trans X_1)>0.
 \label{nc:eq:phase-condition}
\end{equation}
Una familia de estados iniciales se reconstruye mediante
\begin{equation}
 x_{\mathrm{sem}}(\varphi)=X_0+
 2\operatorname{Re}\sum_{k=1}^{H}X_ke^{\ii k\varphi}.
 \label{nc:eq:seed-reconstruction}
\end{equation}

\subsection{Continuación causal}

Defina \(\bar u=\Phi(\bar y)\) y el auxiliar afín
\begin{equation}
 A_{\mathrm a}=A+BKC,\qquad
 d_{\mathrm a}=B(\bar u-K\bar y).
 \label{nc:eq:auxiliary}
\end{equation}
La homotopía exacta es
\begin{equation}
 F_\varepsilon(x)=A_{\mathrm a}x+d_{\mathrm a}
 +\varepsilon B\left[
 \Phi(Cx)-\bar u-K(Cx-\bar y)
 \right],\qquad 0\leq\varepsilon\leq1.
 \label{nc:eq:homotopy}
\end{equation}
Por sustitución directa,
\begin{equation}
 F_0(x)=A_{\mathrm a}x+d_{\mathrm a},
 \qquad F_1(x)=Ax+B\Phi(Cx).
\end{equation}
En Caputo, la continuación debe evaluarse sobre una sola historia:
\begin{equation}
 x(t)=x_{\mathrm{sem}}+
 \frac{1}{\Gamma(q)}\int_{t_0}^{t}
 (t-\tau)^{q-1}F_{\varepsilon(\tau)}(x(\tau))\,d\tau.
 \label{nc:eq:causal-homotopy}
\end{equation}

\subsection{Matignon y frontera de la afirmación de ocultedad}

Para un equilibrio (e) con Jacobiano (J(e)) y espectro
\(\{\lambda_j\}\), el sistema lineal con orden conmensurable \(q\) es
asintóticamente estable si
\begin{equation}
 |\arg\lambda_j|>\frac{q\pi}{2}
 \quad\text{para todo }j.
 \label{nc:eq:matignon}
\end{equation}
Se reportará el margen
\begin{equation}
 m_M(e;q)=\min_j|\arg\lambda_j|-\frac{q\pi}{2}.
 \label{nc:eq:margin}
\end{equation}
Un margen positivo certifica estabilidad lineal local; no certifica caos,
atracción global ni ocultedad.

La ocultedad exige que la cuenca del conjunto atractivo no intersecte
vecindades abiertas de los equilibrios. Un número finito de trayectorias sobre
radios o direcciones discretos sólo sustenta una afirmación bajo el contrato de
muestreo declarado. Esta distinción se conserva en todo el documento.

\clearpage
\section{Caso I: sistema de Lorenz generalizado fraccionario}

\subsection{Modelo, parámetros y simetría}

El sistema registrado es
\begin{equation}
\begin{aligned}
 {}^{\mathrm C}D^qx_1&=-\sigma(x_1-x_2)-a x_2x_3,\\
 {}^{\mathrm C}D^qx_2&=r x_1-x_2-x_1x_3,\\
 {}^{\mathrm C}D^qx_3&=-x_3+x_1x_2,
\end{aligned}
 \qquad \sigma=-ar,\quad a<0.
\label{nc:eq:lorenz-model}
\end{equation}
El caso bibliográfico usa
\begin{equation}
 r=6.8,\qquad a=-0.5,\qquad \sigma=3.4,
 \qquad q=0.995.
 \label{nc:eq:lorenz-parameters}
\end{equation}
El campo es equivariante bajo
\begin{equation}
 \mathcal S(x_1,x_2,x_3)=(-x_1,-x_2,x_3).
 \label{nc:eq:lorenz-symmetry}
\end{equation}
Por tanto, los equilibrios no nulos, semillas y trayectorias aparecen en pares
simétricos.

\subsection{Realización exacta de Lur'e multicanal}

Agrupando términos lineales y monomios se obtiene
\begin{equation}
 A_L=
 \begin{pmatrix}
 -\sigma&\sigma&0\\ r&-1&0\\0&0&-1
 \end{pmatrix},\qquad
 B_L=
 \begin{pmatrix}
 -a&0&0\\0&-1&0\\0&0&1
 \end{pmatrix},
 \qquad C_L=I_3,
 \label{nc:eq:lorenz-abc}
\end{equation}
y
\begin{equation}
 \Phi_L(x)=
 \begin{pmatrix}
 x_2x_3\\x_1x_3\\x_1x_2
 \end{pmatrix}.
 \label{nc:eq:lorenz-phi}
\end{equation}
La comprobación se hace componente a componente:
\begin{align*}
 [A_Lx+B_L\Phi_L(x)]_1
 &=-\sigma x_1+\sigma x_2-a x_2x_3,\\
 [A_Lx+B_L\Phi_L(x)]_2
 &=r x_1-x_2-x_1x_3,\\
 [A_Lx+B_L\Phi_L(x)]_3
 &=-x_3+x_1x_2.
\end{align*}
Luego el residuo de la ecuación~\eqref{nc:eq:lure-residual} es idénticamente
\((0,0,0)\trans\).

\subsection{Derivación completa de la matriz de transferencia}

Escriba \(z=s^q\). Entonces
\begin{equation}
 zI-A_L=
 \begin{pmatrix}
 z+\sigma&-\sigma&0\\
 -r&z+1&0\\
 0&0&z+1
 \end{pmatrix}.
\end{equation}
El determinante del bloque superior es
\begin{equation}
 \Delta_L(z)=(z+\sigma)(z+1)-\sigma r.
 \label{nc:eq:lorenz-delta}
\end{equation}
Por inversión de un bloque \(2\times2\),
\begin{equation}
 (zI-A_L)^{-1}=
 \begin{pmatrix}
 \dfrac{z+1}{\Delta_L}&\dfrac{\sigma}{\Delta_L}&0\\[7pt]
 \dfrac{r}{\Delta_L}&\dfrac{z+\sigma}{\Delta_L}&0\\[7pt]
 0&0&\dfrac{1}{z+1}
 \end{pmatrix}.
\end{equation}
Como \(C_L=I_3\), la multiplicación por \(B_L\) da
\begin{equation}
 W_{L,q}(s)=
 \begin{pmatrix}
 \dfrac{-a(z+1)}{\Delta_L}&
 \dfrac{-\sigma}{\Delta_L}&0\\[8pt]
 \dfrac{-ar}{\Delta_L}&
 \dfrac{-(z+\sigma)}{\Delta_L}&0\\[8pt]
 0&0&\dfrac{1}{z+1}
 \end{pmatrix}_{z=s^q}.
 \label{nc:eq:lorenz-transfer}
\end{equation}
Wolfram redujo a cero las nueve entradas de la diferencia entre
\(C_L(s^qI-A_L)^{-1}B_L\) y~\eqref{nc:eq:lorenz-transfer}.

\subsection{Equilibrios: derivación sin saltos}

En un equilibrio la derivada de Caputo de un estado constante es cero. La
tercera ecuación de~\eqref{nc:eq:lorenz-model} impone
\begin{equation}
 x_3=x_1x_2.
 \label{nc:eq:lorenz-eq1}
\end{equation}
La segunda da
\begin{equation}
 x_2(1+x_1^2)=r x_1,
 \qquad
 x_2=\frac{r x_1}{1+x_1^2}.
 \label{nc:eq:lorenz-eq2}
\end{equation}
Si \(x_1=0\), las ecuaciones restantes producen
\begin{equation}
 E_0=(0,0,0).
\end{equation}
Para \(x_1\neq0\), use \(\sigma=-ar\) en la primera ecuación:
\begin{align}
 0&=ar(x_1-x_2)-a x_2x_3,\\
 r(x_1-x_2)&=x_1x_2^2.
 \label{nc:eq:lorenz-eq3}
\end{align}
Sea \(u=x_1^2\). Al sustituir~\eqref{nc:eq:lorenz-eq2} en
\eqref{nc:eq:lorenz-eq3} se obtiene
\begin{equation}
 (1+u)^2-r(1+2u)=0,
\end{equation}
o, equivalentemente,
\begin{equation}
 u^2+2(1-r)u+(1-r)=0.
\end{equation}
Por la fórmula cuadrática,
\begin{equation}
 u_{\pm}=r-1\pm\sqrt{r(r-1)}.
 \label{nc:eq:lorenz-u}
\end{equation}
Para \(r=6.8\), \(u_-=r-1-\sqrt{r(r-1)}<0\) y no produce equilibrios
reales. Con
\begin{equation}
 u=u_+=r-1+\sqrt{r(r-1)},
\end{equation}
los otros dos equilibrios son
\begin{equation}
 E_{\pm}=\left(
 \pm\sqrt u,
 \pm\frac{r\sqrt u}{1+u},
 \frac{ru}{1+u}
 \right).
 \label{nc:eq:lorenz-equilibria}
\end{equation}
Numéricamente,
\begin{equation}
 E_{\pm}=(\pm3.47564776512854,
 \ \pm1.80689408468029,
 \ 6.28012738724303).
 \label{nc:eq:lorenz-equilibria-num}
\end{equation}
La sustitución simbólica de~\eqref{nc:eq:lorenz-equilibria} en el campo produjo
residuo exactamente nulo; el mayor residuo numérico de los tres equilibrios fue
\(1.07\times10^{-14}\).

\subsection{Jacobiano, polinomios y corrección espectral}

El Jacobiano correcto es
\begin{equation}
 J_L(x)=
 \begin{pmatrix}
 -\sigma&\sigma-a x_3&-a x_2\\
 r-x_3&-1&-x_1\\
 x_2&x_1&-1
 \end{pmatrix}.
 \label{nc:eq:lorenz-jacobian}
\end{equation}
En el origen,
\begin{equation}
 p_0(\lambda)
 =(\lambda+1)\left[(\lambda+\sigma)(\lambda+1)-\sigma r\right].
\end{equation}
Con~\eqref{nc:eq:lorenz-parameters},
\begin{equation}
 p_0(\lambda)=
 \lambda^3+5.4\lambda^2-15.32\lambda-19.72,
\end{equation}
y
\begin{equation}
 \operatorname{spec}J_L(E_0)=
 \{-7.15580467734555,\ 2.75580467734555,\ -1\}.
 \label{nc:eq:lorenz-origin-spectrum}
\end{equation}
El valor propio real positivo hace inestable a \(E_0\) para todo \(q>0\).

En cualquiera de los equilibrios \(E_{\pm}\), el polinomio mónico es
\begin{equation}
 p_{\pm}(\lambda)=
 \lambda^3+5.4\lambda^2
 +14.8476942706167\lambda+78.88,
\end{equation}
con espectro
\begin{equation}
 \operatorname{spec}J_L(E_{\pm})=
 \left\{
 -5.37029951773135,
 -0.0148502411343235\pm3.83248917224189\ii
 \right\}.
 \label{nc:eq:lorenz-outer-spectrum}
\end{equation}
Para el par complejo,
\begin{equation}
 q_{\mathrm{crit}}=
 \frac{2|\arg\lambda|}{\pi}=1.00246678056773,
\end{equation}
de modo que \(E_{\pm}\) es localmente estable para \(q=0.995\). El margen es
\begin{equation}
 m_M(E_{\pm};0.995)=0.0117287914887680>0.
\end{equation}

\begin{hallazgo}{Corrección algebraica al dato bibliográfico heredado}
La fuente primaria y la extracción JSON local registran en \(E_0\) el espectro
\(\{2.5576,-1,-7.5576\}\) y en \(E_{\pm}\) el espectro
\(\{-5.9570,-0.0215\pm3.6026\ii\}\). Esos números no corresponden a las
ecuaciones y parámetros declarados: ambos conjuntos suman \(-6\), mientras que
\(\operatorname{tr}J_L=-5.4\). Wolfram dio residuos máximos \(27.17\) y
\(29.33\) al sustituirlos en los polinomios característicos. Los valores
corregidos son~\eqref{nc:eq:lorenz-origin-spectrum} y
\eqref{nc:eq:lorenz-outer-spectrum}. La clasificación cualitativa no cambia:
origen inestable y par externo estable. También debe evitarse llamar
``silla'' a un equilibrio declarado asintóticamente estable.
\end{hallazgo}

\subsection{Sistema algebraico de la semilla armónica}

Para \(H\geq2\), escriba \(X_{j,k}\) para el coeficiente \(k\) de \(x_j\).
Los coeficientes no lineales son
\begin{equation}
 \widehat\Phi_{L,k}=
 \begin{pmatrix}
 \displaystyle\sum_m X_{2,m}X_{3,k-m}\\[5pt]
 \displaystyle\sum_m X_{1,m}X_{3,k-m}\\[5pt]
 \displaystyle\sum_m X_{1,m}X_{2,k-m}
 \end{pmatrix}.
 \label{nc:eq:lorenz-fourier-phi}
\end{equation}
Por tanto, una semilla no se obtiene escogiendo una condición inicial al azar,
sino resolviendo
\begin{align}
 A_LX_0+B_L\widehat\Phi_{L,0}&=0,\\
 \left[(\ii k\omega)^qI-A_L\right]X_k
 -B_L\widehat\Phi_{L,k}&=0,
 \quad k=1,\ldots,H,\\
 \operatorname{Im}(e_j\trans X_1)&=0.
 \label{nc:eq:lorenz-hb-system}
\end{align}
La linealización alrededor de un sesgo \(\bar x\) usa
\begin{equation}
 K_L(\bar x)=D\Phi_L(\bar x)=
 \begin{pmatrix}
 0&\bar x_3&\bar x_2\\
 \bar x_3&0&\bar x_1\\
 \bar x_2&\bar x_1&0
 \end{pmatrix}
 \label{nc:eq:lorenz-k}
\end{equation}
y el predictor local
\begin{equation}
 \det\left[I_3-W_{L,q}(\ii\omega)K_L(\bar x)\right]=0.
 \label{nc:eq:lorenz-local-closure}
\end{equation}
La solución de~\eqref{nc:eq:lorenz-hb-system} debe verificarse con \(H\) y
\(H+1\), reconstruirse mediante~\eqref{nc:eq:seed-reconstruction} y transportarse
con~\eqref{nc:eq:causal-homotopy}.

\subsection{Datos y comprobaciones necesarios para una búsqueda local}

La referencia declara \(q=0.995\), un atractor representado gráficamente y 50
trayectorias desde una vecindad de radio \(0.3\) del origen que llegan a los dos
equilibrios estables. Sin embargo, la evidencia disponible no especifica la
condición inicial del atractor, el paso, el horizonte, el transitorio, la
historia ni una semilla de balance armónico. Para cerrar esa laguna se requiere:
\begin{enumerate}
 \item resolver~\eqref{nc:eq:lorenz-hb-system} para \(H=2,3,\ldots\) y retener
 sólo raíces con residuos y distorsión convergentes;
 \item reconstruir todas las fases no equivalentes por
 \eqref{nc:eq:seed-reconstruction};
 \item continuar en una sola historia de Caputo;
 \item reiniciar una evaluación autónoma separada desde el estado terminal;
 \item demostrar atracción local y robustez frente a (h), memoria, horizonte
 y solver;
 \item sondear todos los equilibrios con bolas y superficies declaradas, no
 únicamente 50 direcciones alrededor de \(E_0\).
\end{enumerate}

\clearpage
\section{Caso II: sistema de Rabinovich--Fabrikant fraccionario}

\subsection{Modelo, parámetros y simetría}

El sistema RF registrado es
\begin{equation}
\begin{aligned}
 {}^{\mathrm C}D^qx_1&=x_2(x_3-1+x_1^2)+a x_1,\\
 {}^{\mathrm C}D^qx_2&=x_1(3x_3+1-x_1^2)+a x_2,\\
 {}^{\mathrm C}D^qx_3&=-2x_3(b+x_1x_2),
\end{aligned}
\label{nc:eq:rf-model}
\end{equation}
con
\begin{equation}
 a=0.1,\qquad b=0.2876,\qquad q=0.998.
 \label{nc:eq:rf-parameters}
\end{equation}
También es equivariante bajo
\begin{equation}
 \mathcal S(x_1,x_2,x_3)=(-x_1,-x_2,x_3).
\end{equation}

\subsection{Realización polinómica multicanal exacta}

Expanda el campo:
\begin{align*}
 F_1&=a x_1-x_2+x_2x_3+x_1^2x_2,\\
 F_2&=x_1+a x_2+3x_1x_3-x_1^3,\\
 F_3&=-2b x_3-2x_1x_2x_3.
\end{align*}
Una realización directa, sin variables auxiliares dinámicas, es
\begin{equation}
 A_R=
 \begin{pmatrix}
 a&-1&0\\1&a&0\\0&0&-2b
 \end{pmatrix},\qquad C_R=I_3,
 \label{nc:eq:rf-a}
\end{equation}
\begin{equation}
 B_R=
 \begin{pmatrix}
 1&1&0&0&0\\
 0&0&3&-1&0\\
 0&0&0&0&-2
 \end{pmatrix},
 \label{nc:eq:rf-b}
\end{equation}
y
\begin{equation}
 \Phi_R(x)=
 \begin{pmatrix}
 x_2x_3\\x_1^2x_2\\x_1x_3\\x_1^3\\x_1x_2x_3
 \end{pmatrix}.
 \label{nc:eq:rf-phi}
\end{equation}
En efecto,
\begin{equation}
 A_Rx+B_R\Phi_R(x)=F_R(x)
\end{equation}
componente a componente. Wolfram simplificó el residuo simbólico a
\((0,0,0)\trans\).

Si una implementación exige potencias de mediciones lineales, los productos
pueden factorizarse sin aproximación mediante
\begin{equation}
 uv=\frac{(u+v)^2-(u-v)^2}{4}
 \label{nc:eq:square-identity}
\end{equation}
y
\begin{equation}
 uvw=\frac{(u+v+w)^3+(u-v-w)^3+(-u+v-w)^3+(-u-v+w)^3}{24}.
 \label{nc:eq:cubic-identity}
\end{equation}
Estas identidades cambian el número de canales, no el campo vectorial.

\subsection{Derivación completa de la transferencia RF}

De nuevo, sea \(z=s^q\). Entonces
\begin{equation}
 zI-A_R=
 \begin{pmatrix}
 z-a&1&0\\-1&z-a&0\\0&0&z+2b
 \end{pmatrix}.
\end{equation}
Defina
\begin{equation}
 \Delta_R(z)=(z-a)^2+1.
 \label{nc:eq:rf-delta}
\end{equation}
La inversa es
\begin{equation}
 (zI-A_R)^{-1}=
 \begin{pmatrix}
 \dfrac{z-a}{\Delta_R}&\dfrac{-1}{\Delta_R}&0\\[7pt]
 \dfrac{1}{\Delta_R}&\dfrac{z-a}{\Delta_R}&0\\[7pt]
 0&0&\dfrac{1}{z+2b}
 \end{pmatrix}.
\end{equation}
Al multiplicar por las cinco columnas de \(B_R\),
\begin{equation}
 W_{R,q}(s)=
 \begin{pmatrix}
 \dfrac{z-a}{\Delta_R}&
 \dfrac{z-a}{\Delta_R}&
 \dfrac{-3}{\Delta_R}&
 \dfrac{1}{\Delta_R}&0\\[8pt]
 \dfrac{1}{\Delta_R}&
 \dfrac{1}{\Delta_R}&
 \dfrac{3(z-a)}{\Delta_R}&
 \dfrac{-(z-a)}{\Delta_R}&0\\[8pt]
 0&0&0&0&\dfrac{-2}{z+2b}
 \end{pmatrix}_{z=s^q}.
 \label{nc:eq:rf-transfer}
\end{equation}
La diferencia simbólica entre~\eqref{nc:eq:rf-transfer} y
\((s^qI-A_R)^{-1}B_R\) fue la matriz nula \(3\times5\).

\subsection{Equilibrios: reducción a una ecuación cuadrática}

El origen
\begin{equation}
 E_0=(0,0,0)
\end{equation}
es inmediato. Para un equilibrio con \(x_3\neq0\), la tercera ecuación exige
\begin{equation}
 x_1x_2=-b,\qquad x_2=-\frac{b}{x_1}.
 \label{nc:eq:rf-product}
\end{equation}
La primera ecuación, multiplicada por \(x_1\), da
\begin{equation}
 -b(x_3-1+x_1^2)+a x_1^2=0,
\end{equation}
de donde
\begin{equation}
 x_3=1+\frac{a-b}{b}x_1^2.
 \label{nc:eq:rf-z}
\end{equation}
Sea \(u=x_1^2\). La segunda ecuación, multiplicada por \(x_1\), es
\begin{equation}
 u(3x_3+1-u)-ab=0.
\end{equation}
Sustituyendo~\eqref{nc:eq:rf-z}, se llega a
\begin{equation}
 (3a-4b)u^2+4bu-ab^2=0.
 \label{nc:eq:rf-u-polynomial}
\end{equation}
Con
\begin{equation}
 R=\sqrt{4+a(3a-4b)},
\end{equation}
las dos raíces son
\begin{equation}
 u_{\pm}=\frac{b(-2\pm R)}{3a-4b}.
 \label{nc:eq:rf-u}
\end{equation}
Para~\eqref{nc:eq:rf-parameters}, ambas son positivas. Cada una genera un par
simétrico:
\begin{equation}
 E_{\pm}^{(u)}=\left(
 \pm\sqrt u,\ \mp\frac{b}{\sqrt u},
 1+\frac{a-b}{b}u
 \right).
 \label{nc:eq:rf-equilibria-general}
\end{equation}
Ordenando por \(|x_1|\), se obtienen
\begin{align}
 E_{1,2}&=(\mp1.15997695583801,
 \ \pm0.247935959893466,
 \ 0.122306917444673),
 \label{nc:eq:rf-e12}\\
 E_{3,4}&=(\mp0.085021329989515,
 \ \pm3.38268055834300,
 \ 0.995284804098130).
 \label{nc:eq:rf-e34}
\end{align}
Wolfram redujo exactamente a cero el campo evaluado en las dos ramas
algebraicas de~\eqref{nc:eq:rf-u}; el máximo residuo numérico fue
\(1.06\times10^{-14}\).

\subsection{Jacobiano, polinomios y Matignon}

La diferenciación de~\eqref{nc:eq:rf-model} da
\begin{equation}
 J_R(x)=
 \begin{pmatrix}
 2x_1x_2+a&x_3-1+x_1^2&x_2\\
 3x_3+1-3x_1^2&a&3x_1\\
 -2x_2x_3&-2x_1x_3&-2(b+x_1x_2)
 \end{pmatrix}.
 \label{nc:eq:rf-jacobian}
\end{equation}
En el origen,
\begin{equation}
 p_0(\lambda)=(\lambda+2b)\left[(\lambda-a)^2+1\right],
\end{equation}
y
\begin{equation}
 \operatorname{spec}J_R(E_0)=
 \{-0.5752,\ 0.1\pm1.0\ii\}.
\end{equation}
El par complejo tiene
\begin{equation}
 q_{\mathrm{crit}}=0.936548965138893.
\end{equation}
Como \(q=0.998>q_{\mathrm{crit}}\), \(E_0\) es inestable y
\begin{equation}
 m_M(E_0;0.998)=-0.0965270598375723.
\end{equation}

En \(E_{1,2}\),
\begin{equation}
 p_{1,2}(\lambda)=
 \lambda^3+0.3752\lambda^2
 +2.20397205132125\lambda+0.556792727146772,
\end{equation}
\begin{equation}
 \operatorname{spec}J_R(E_{1,2})=
 \{-0.256175578455019,
 -0.059512210772491\pm1.47307134859472\ii\}.
\end{equation}
El mínimo crítico es
\begin{equation}
 q_{\mathrm{crit}}=1.02570551509697>1,
\end{equation}
por lo que el par \(E_{1,2}\) es estable para todo \(0<q\leq1\); en particular,
\begin{equation}
 m_M(E_{1,2};0.998)=0.0435197213462810>0.
\end{equation}

En \(E_{3,4}\),
\begin{equation}
 p_{3,4}(\lambda)=
 \lambda^3+0.3752\lambda^2
 +22.7628315898639\lambda-4.53095664529540,
\end{equation}
\begin{equation}
 \operatorname{spec}J_R(E_{3,4})=
 \{0.198062714510438,
 -0.286631357255219\pm4.77432885589854\ii\}.
\end{equation}
El valor propio positivo implica inestabilidad para todo \(q>0\). Estos
resultados coinciden, dentro del redondeo impreso, con la fuente primaria.

\subsection{Sistema algebraico de la semilla RF}

El Jacobiano estático de~\eqref{nc:eq:rf-phi}, evaluado en un sesgo \(\bar x\), es
\begin{equation}
 K_R(\bar x)=
 \begin{pmatrix}
 0&\bar x_3&\bar x_2\\
 2\bar x_1\bar x_2&\bar x_1^2&0\\
 \bar x_3&0&\bar x_1\\
 3\bar x_1^2&0&0\\
 \bar x_2\bar x_3&\bar x_1\bar x_3&\bar x_1\bar x_2
 \end{pmatrix}.
 \label{nc:eq:rf-k}
\end{equation}
El predictor local satisface
\begin{equation}
 \det\left[I_3-W_{R,q}(\ii\omega)K_R(\bar x)\right]=0.
 \label{nc:eq:rf-local-closure}
\end{equation}

Para el balance de amplitud finita se recomienda iniciar con \(H=3\). Las cinco
componentes no lineales se calculan sin ambigüedad por
\begin{equation}
 \widehat\Phi_{R,k}=
 \begin{pmatrix}
 \displaystyle\sum_m X_{2,m}X_{3,k-m}\\[5pt]
 \displaystyle\sum_{m,n}X_{1,m}X_{1,n}X_{2,k-m-n}\\[5pt]
 \displaystyle\sum_m X_{1,m}X_{3,k-m}\\[5pt]
 \displaystyle\sum_{m,n}X_{1,m}X_{1,n}X_{1,k-m-n}\\[5pt]
 \displaystyle\sum_{m,n}X_{1,m}X_{2,n}X_{3,k-m-n}
 \end{pmatrix}.
 \label{nc:eq:rf-fourier-phi}
\end{equation}
El sistema de semillas es
\begin{align}
 A_RX_0+B_R\widehat\Phi_{R,0}&=0,\\
 \left[(\ii k\omega)^qI-A_R\right]X_k
 -B_R\widehat\Phi_{R,k}&=0,
 \quad k=1,\ldots,H,\\
 \operatorname{Im}(e_j\trans X_1)&=0.
 \label{nc:eq:rf-hb-system}
\end{align}
La contribución cúbica genera armónicos por encima de \(H\). Por ello, cada raíz
debe continuar a \(H+1\), y la cola omitida debe registrarse como índice de
distorsión. Sólo después se reconstruye~\eqref{nc:eq:seed-reconstruction} y se
aplica la homotopía~\eqref{nc:eq:homotopy}.

\subsection{Qué falta para una búsqueda local reproducible}

La fuente usa \(q=0.998\) y vecindades de radio \(0.1\) alrededor de los
equilibrios inestables \(E_0,E_3,E_4\). La evidencia disponible no incluye la
condición inicial publicada del conjunto verde, una semilla algebraica, el
contrato ABM completo ni las trayectorias de vecindad. Para cerrar la búsqueda
se requiere:
\begin{enumerate}
 \item resolver~\eqref{nc:eq:rf-hb-system} desde varias ramas y fases;
 \item separar soluciones acotadas, divergentes y transitoriamente caóticas;
 \item continuar cada semilla con una única historia causal;
 \item comprobar la dinámica con un segundo integrador de memoria completa;
 \item sondear los cinco equilibrios y registrar explícitamente los destinos
 \(E_{1,2}\), divergencia, conjunto objetivo u otros conjuntos;
 \item distinguir una ausencia finita de contactos de una prueba de vecindad
 abierta.
\end{enumerate}

\section{Comparación y prioridad de implementación}

\begin{table}[ht]
\centering
\small
\caption{Comparación algebraica de los dos candidatos prioritarios.}
\begin{tabularx}{\linewidth}{>{\raggedright\arraybackslash}p{0.24\linewidth}XX}
\toprule
Propiedad & Lorenz generalizado & Rabinovich--Fabrikant\\
\midrule
Orden no lineal & cuadrático & cuadrático y cúbico\\
Canales elegidos & 3 & 5\\
Truncación inicial & \(H=2\) & \(H=3\)\\
Equilibrios & 3 & 5\\
Estables a \(q\) objetivo & \(E_+,E_-\) & \(E_1,E_2\)\\
Inestables a \(q\) objetivo & \(E_0\) & \(E_0,E_3,E_4\)\\
Simetría & \((x_1,x_2,x_3)\mapsto(-x_1,-x_2,x_3)\) & la misma\\
Riesgo algebraico & bajo--medio & medio--alto\\
Riesgo numérico & historia y cuenca & historia, cubics, divergencia y transitorios\\
Estado local & referencia solamente & referencia solamente\\
\bottomrule
\end{tabularx}
\end{table}

La prioridad recomendada es implementar primero Lorenz generalizado. Su campo
es cuadrático, requiere menos canales y \(H=2\), y la cuenca publicada alrededor
del único equilibrio inestable ofrece un control inicial simple. RF debe ser el
segundo caso: es científicamente valioso, pero los términos cúbicos, los tres
equilibrios inestables y las trayectorias divergentes elevan el costo de una
decisión fiable.

Esta prioridad no es una afirmación sobre cuál atractor es ``más oculto''. Es
una decisión de verificabilidad: conviene cerrar primero el caso con menor
número de grados armónicos y de destinos de cuenca.

\section{Verificación algebraica cruzada con Wolfram Language}

Las comprobaciones se ejecutaron el 10 de agosto de 2026 en un kernel
Wolfram Language sin estado persistente. Se usaron parámetros racionales
\(r=34/5\), \(a_L=-1/2\), \(a_R=1/10\) y \(b_R=719/2500\) para evitar que el
redondeo antecediera al álgebra.

\begin{longtable}{>{\raggedright\arraybackslash}p{0.34\linewidth}
>{\raggedright\arraybackslash}p{0.17\linewidth}
>{\raggedright\arraybackslash}p{0.40\linewidth}}
\caption{Matriz de verificación Wolfram.}\\
\toprule
Comprobación & Resultado & Evidencia\\
\midrule
\endfirsthead
\toprule
Comprobación & Resultado & Evidencia\\
\midrule
\endhead
Identidad Lur'e de Lorenz & \ok &
\(A_Lx+B_L\Phi_L(x)-F_L(x)=0\) en las tres componentes.\\
Transferencia de Lorenz & \ok &
Las nueve entradas del residuo simbólico son cero.\\
Equilibrios de Lorenz & \ok &
Residuo simbólico cero; máximo numérico \(1.07\times10^{-14}\).\\
Espectros impresos de Lorenz & \corr &
No satisfacen los polinomios declarados; la traza y los residuos revelan la
inconsistencia. Se proporcionan espectros corregidos.\\
Clasificación de Matignon de Lorenz & \ok &
\(E_0\) inestable; \(E_{\pm}\) estables a \(q=0.995\), margen \(0.0117288\).\\
Identidad Lur'e de RF & \ok &
\(A_Rx+B_R\Phi_R(x)-F_R(x)=0\) en las tres componentes.\\
Transferencia de RF & \ok &
Las quince entradas del residuo simbólico son cero.\\
Equilibrios RF & \ok &
Las dos ramas de~\eqref{nc:eq:rf-u} dan residuo simbólico cero; máximo numérico
\(1.06\times10^{-14}\).\\
Espectros y Matignon de RF & \ok &
Coinciden con la fuente al redondeo: \(E_{1,2}\) estables;
\(E_0,E_{3,4}\) inestables a \(q=0.998\).\\
\bottomrule
\end{longtable}

La verificación Wolfram cubre la realización algebraica. No se ejecutó una
integración de Caputo ni una campaña de cuencas en Wolfram; por tanto, este
resultado no debe etiquetarse como reproducción dinámica.

\section{Conclusiones}

\begin{enumerate}
 \item Los dos candidatos fraccionarios no Chua con procedencia directa de
 atractor oculto son Lorenz generalizado y Rabinovich--Fabrikant.
 \item Ambos admiten una realización de Lur'e multicanal exacta. Para Lorenz se
 obtuvo una transferencia \(3\times3\); para RF, una transferencia \(3\times5\).
 \item Se derivaron todos los equilibrios desde las ecuaciones, sus Jacobianos,
 polinomios característicos y márgenes de Matignon.
 \item Se detectó y corrigió una inconsistencia numérica en los espectros de
 Lorenz heredados de la fuente primaria. La clasificación cualitativa no cambia.
 \item Las ecuaciones~\eqref{nc:eq:lorenz-hb-system} y
 \eqref{nc:eq:rf-hb-system} son la realización algebraica completa necesaria para
 generar semillas; todavía no hay raíces numéricas promovibles.
 \item La próxima evidencia científica debe ser dinámica: solución armónica
 convergente, continuación causal, evaluación autónoma, robustez y sondeos de
 todos los equilibrios. Hasta entonces, ambos casos se conservan como evidencia
 bibliográfica y algebraica, sin validación dinámica.
\end{enumerate}
